\section{Introduction}

A self-replicating machine that builds copies of itself from local
material has been a fixture of long-range space-systems thinking since von Neumann's
self-reproducing automata~\citep{von-neumann-burks-1966} and Freitas' self-reproducing
interstellar probe~\citep{freitas-1980-interstellar-probe}. Its appeal is leverage: a
seed that remakes most of its own mass in place turns each launched kilogram into far
more installed capacity than a finished payload could~\citep{freitas-merkle-2004}. The
same leverage, applied recursively, is what makes a single probe a plausible way to
cross a galaxy, and even to spread between galaxies, in a time short against its
age~\citep{armstrong-sandberg-2013}: the observation that sharpens the Fermi
question~\citep{hart-1975,tipler-1980}, though its force is itself contested~\citep{sagan-newman-1983}.

That programme is naturally studied at three very different scales, and in practice each
is modelled on its own. At the smallest scale a \emph{single factory} lands, closes on
as much of its bill of materials as it can, and copies itself. At an intermediate scale a
\emph{local fleet} of such factories disperses across a solar system, each child building
the next. At the largest scale a \emph{settlement front} of self-replicating probes
crosses interstellar space, and the exploration timescale is set by how the front chooses
targets and how fast the probes travel~\citep{nicholson-forgan-2013,forgan-papadog-kitching-2013}.

Modelling the three separately leaves a seam. Every scale needs one shared number, the
time to build a single copy, and each has been free to pick its own. In the settlement-front
model we build on, that number appears as a per-star \emph{manufacturing dwell}: the time a
freshly settled probe spends building offspring before they depart. For want of a sourced
value it has been set to zero in that model, so that the front replicates instantaneously and
only travel costs time, an idealization it inherits from its settlement-front
predecessors~\citep{nicholson-forgan-2013,forgan-papadog-kitching-2013}, none of which derives
a build time. Whether that zero is safe has not been derived from the factory physics that
actually fixes the dwell; the nearest prior work sweeps a dwell-like imprint time for
detectability rather than for its cost on the fill~\citep{cotta-morales-2009}.

This paper closes the seam and asks what the closed form says. We thread one sourced factory
through all three folds and \emph{derive} the manufacturing dwell from the very build
physics the fleet already uses, so the number is computed once, from cited inputs, rather
than assumed three times over. We introduce no new physical quantity; we route one the
factory already fixes to the scale that was guessing it. The question we can then answer is
cross-scale and quantitative: \emph{which constraint binds at which scale?} The answer,
developed below, is that the manufacturing cadence that is the binding clock for a local
fleet is a negligible tax on a powered interstellar cruise, robust to a hundredfold uncertainty
in the one number it rests on, and a small but rising tax as the front is driven to slingshot
or directed-energy speeds.
